\documentclass[12pt]{article}
\usepackage[english, bulgarian]{babel}
\usepackage[utf8]{inputenc}
\usepackage[T2A]{fontenc}
\usepackage{amssymb}
\usepackage{hyperref, fancyhdr, lastpage, fancyvrb, tcolorbox, titlesec}
\usepackage{array, tabularx, colortbl}
\usepackage{tikz}
\usepackage{venndiagram}
\usepackage{amsthm, bm}
\usepackage{relsize}
\usepackage{amsmath,physics}
\usepackage{mathtools}
\usepackage{subcaption}
\usepackage{scalerel}
\usepackage{theoremref}
\usepackage{circuitikz}
\usepackage{geometry}
\usepackage{stmaryrd}
\usepackage{forest}
\usepackage{cancel}
\usepackage{faktor}
\usepackage{gensymb}
\usetikzlibrary{automata, arrows, positioning, shapes}
\usetikzlibrary{decorations.pathmorphing}
\useforestlibrary{linguistics}

\ExplSyntaxOn
\NewDocumentCommand{\opair}{m}
 {
  \langle\mspace{2mu}
  \clist_set:Nn \l_tmpa_clist { #1 }
  \clist_use:Nn \l_tmpa_clist {,\mspace{3mu plus 1mu minus 1mu}\allowbreak}
  \mspace{2mu}\rangle
}
\ExplSyntaxOff

\hypersetup{
    colorlinks=true,
    linktoc=all,
    linkcolor=blue
}

\setlength\parindent{0pt}

\newcommand{\N}{\mathbb{N}}
\newcommand{\Z}{\mathbb{Z}}
\newcommand{\Q}{\mathbb{Q}}
\newcommand{\R}{\mathbb{R}}
\newcommand{\E}{\mathbb{E}}

\newcommand{\vars}{\operatorname{Vars}}
\newcommand{\free}{\operatorname{Free}}

\newcommand{\logand}{\; \& \;}

\newcommand{\calA}{\mathcal{A}}
\newcommand{\calS}{\mathcal{S}}
\newcommand{\calD}{\mathcal{D}}
\newcommand{\calL}{\mathcal{L}}
\newcommand{\calZ}{\mathcal{Z}}
\newcommand{\calQ}{\mathcal{Q}}
\newcommand{\calR}{\mathcal{R}}
\newcommand{\calN}{\mathcal{N}}
\newcommand{\calM}{\mathcal{M}}
\newcommand{\calF}{\mathcal{F}}
\newcommand{\calP}{\mathcal{P}}
\newcommand{\calC}{\mathcal{C}}
\newcommand{\calG}{\mathcal{G}}
\newcommand{\calB}{\mathcal{B}}

\newcommand{\fa}{\kern.05em\ensuremath\forall\kern-.22em\rotatebox{69}{\rule{.73em}{.4pt}}}
\newcommand{\ex}{\kern.05em\ensuremath\exists\kern-.03em\rotatebox{90}{\rule{.68em}{.35pt}}}


\newcommand{\dequiv}{\stackrel{\text{деф.}}{\longleftrightarrow}}

\newcommand{\db}[1]{\llbracket #1 \rrbracket}

\DeclareMathOperator*{\bigand}{\scalerel*{\&}{\prod}}
\DeclareMathOperator*{\bigor}{\scalerel*{\lor}{\sum}}

\newtheorem*{definition}{Дефиниция}
\newtheorem*{claim}{Твърдение}
\newtheorem*{property}{Свойство}
\newtheorem*{hint}{Упътване}
\theoremstyle{definition}
\newtheorem*{solution}{Решение}
\newtheorem*{notation}{Нотация}
\newtheorem{problem}{Задача}

\title{Рубрика `Често срещани грешки по ЛП`}
% \author{}
\date{}

\begin{document}
\maketitle

\section{Смесване на обектния език с метаезика}

\subsection*{Кои са двата езика?}
\textbf{Обектният език} е предикатния език в задачата, който сме фиксирали.
В него сме фиксирали:
\begin{itemize}
  \item променливи (които по конвенция бележим с $x, y, z$ и потенциални украшения);
  \item константни символи (които по конвенция бележим с $a, b, c$ и потенциални украшения);
  \item функционални символи (които по конвенция бележим с $f, g$;
  \item предикатни символи (които по конвенция бележим с $p, q, r$);
  \item символите $\forall, \exists, \doteq, \Rightarrow, \Leftrightarrow, \neg, \lor, \&, (, )$ и $,$ (символа за запетая).
\end{itemize}
Формулите, които пишем, са някакви крайни низове, които използват \textbf{САМО} тези символи.
Това са редици от символи, които нямат никакво значение, и без намеса от метаезика не носят \textbf{НИКАКВО} значение.
Хората понякога пишат формули, които съдържат неща от вида:
\[
  p^\calM(x, y) \text{ или } f^\calS(z, x).
\]
Горните неща не са конструкции на обектния език.
По нашите конвенции $p^\calM$ е релация (или функция, която връща И/Л) и $f^\calS$ е функция.
Изразът $p^\calM(x, y)$ е безсмислен, тъй като в общия случай $x, y$ са променливи, а не елементи от $|\calM|$.
Същото може да се каже и за $f^\calS(z, x)$, където в общия случай променливите $z$ и $x$ не са елементи на $|\calS|$.

Друга често срещана грешка е добавянето на нова константа в езика, например се наблюдава формула от вида:
\[
  \tau \doteq 0 \text{ или } \tau \doteq a.
\]
Човек вижда, че $0$ или $a$ e от универсума, и директно решава, че ще ги използва.
Забелязвам, че доста често се използва константата $0$, за да се определи $\{ 0 \}$.
Първо, че тя не е част от \textbf{ОБЕКТНИЯ ЕЗИК}.
Второ, дори и да беше част от него, тогава щеше напълно да се обезсмисли това да се определя самият обект, като той може да бъде достъпен със съответния символ за константа.

\textbf{Метаезикът} е езикът, на който правим математически разсъждения (български + някои математически символи, като $=, \in$).
Тук \textbf{ВСИЧКО}, което пишем, трябва да \textbf{ИМА} значение.
Хората, когато доказват, че $f^\calA$ (интерпретацията на функционален символ $f$ в структура $\calA$) е асоциативна, понякога пишат текст от вида:
\[
  f(x, f(y, z)) = f(f(x, y), z) \text{ за всяко } x, y, z.
\]
Това, което горе пише, е че два терма (редици от символи от обектния език) са равни (буква по буква), което не е вярно.
Правилното нещо, което трябва да се напише в случая е:
\[
  f^\calA(a_1, f^\calA(a_2, a_3) = f^\calA(f^\calA(a_1, a_2), a_3) \text{ за всяко } a_1, a_2, a_3 \in |\calA|.
\]
Тук вече няма термове, а има някакви стойности от универсума.
Казвам, че прилагайки функцията върху тези аргументи по тези два начина, ще получа един и същ резултат, без значение какви аргументи взема от универсума, което се и иска от нас.
Някои пък са писали и неща от вида:
\[
  f(x, f(y, z)) = f(f(x, y), z) \text{ за } \forall x \forall y \forall z.
\]
Първо, в този курс сам по себе си символът $\forall$ не носи никакво значение.
Трябва вместо $\forall$ да се пише \textit{за всяко}.
Ако някой толкова го мързи да пише този текст, може да използва:
\begin{itemize}
  \item $\fa$, като съкращение (на \textbf{МЕТАЕЗИКА}) на \textit{за всяко};
  \item $\ex$, като съкращение (на \textbf{МЕТАЕЗИКА}) на \textit{съществува}.
\end{itemize}

\newpage

\subsection*{Как се доказва, че $X \subseteq |\calA|^n$ е определимо?}
\begin{enumerate}
  \item Пишем \textbf{ФОРМУЛА} на \textbf{ДАДЕНИЯ В УСЛОВИЕТО} предикатен език $\varphi$.
  \item Пишем:
    \begin{align*}
      \calA \models \varphi \db{a_1, \dots, a_n} & \longleftrightarrow \text{ превод на `} \calA \models \varphi \text{` на метаезика по дефиницията } \\
                                                 & \stackrel{(\star)}{\longleftrightarrow} \opair{a_1, \dots, a_n} \in X.
    \end{align*}
  \item Доказваме $(\star)$, което става в \textbf{МЕТАЕЗИКА}.
\end{enumerate}

\subsection*{Как се доказва, че $X \subseteq |\calA|^n$ не е определимо?}
\begin{enumerate}
  \item Построяваме (в \textbf{МЕТАЕЗИКА}) автоморфизъм $h$, за който $h(X) \neq X$.
  \item Доказваме (в \textbf{МЕТАЕЗИКА}), че $h : |\calA| \rightarrow |\calA|$ е биекция.
  \item Доказваме (в \textbf{МЕТАЕЗИКА}), че $h$ е силен хомоморфизъм.
  \item Доказваме (в \textbf{МЕТАЕЗИКА}), че $h(X) \neq X$.
\end{enumerate}

\subsection*{Как се доказва, че едно множество от затворени формули $\Gamma$ е изпълнимо?}
\begin{enumerate}
  \item Посочваме структура $\calA$ за \textbf{ОБЕКТНИЯ ЕЗИК} (като посочването става на ниво \textbf{МЕТАЕЗИК}).
  \item Доказваме (в \textbf{МЕТАЕЗИКА}), че $\calA \models \Gamma$, тоест за всяко $\varphi \in \Gamma$ е изпълнено $\calA \models \varphi$.
\end{enumerate}

\subsection*{Защо не е хубаво да се смесват двата езика?}
Представете си, че трябва да направите речник, който превежда български думи на английски.
Това е аналог на задачите за определимост.
Майсторлъка в писането на речник се състои в избора на едно малко множество от прости думи на английски, което да е достатъчно за превод на възможно най-много български думи.
В задачата за определимост предикатния език и структурата играят роля на английския, тоест езика, който се предполага, че читателят знае.
Двойката от множеството, което определяме, и формулата, която го определя, е статията в речника.

\newpage

Следните неща биха направили човек, който си закупи нашия речник, много недоволен:
\begin{itemize}
  \item Той потърси значението на някоя дума, и в нейната дефиниция срещне символ, който не е част от английската азбука, например:
        \begin{center}
          \textit{диван -- canapé} или \textit{стол -- стол}.
        \end{center}
        Това съответства на глупости от вида (тук $0$ не е в езика):
        \[
          \varphi_1[x] \leftrightharpoons x \doteq s(0) \text{ или } \varphi_0[x] \leftrightharpoons x \doteq 0.
        \]
  \item Той потърси значението на някоя дума, и в нейната дефиниция има сложни думи на английски, които никой не използва, например:
        \begin{center}
          \textit{безполезен -- something that can be floccinaucinihilipilificated}
        \end{center}
        Това съответства на глупости от вида:
        \[
          \varphi_1[x] \leftrightharpoons (\forall y \in \N) (\varphi_<[y, z] \Rightarrow \varphi_0[y]).
        \]
        Ограничения на квантори никой не е вкарвал в нашия език, значи не можем да си позволяваме да го правим.
        Правилното нещо би било:
        \[
          \varphi_1[x] \leftrightharpoons \forall y (\varphi_\N[y] \logand \varphi_<[y, z] \Rightarrow \varphi_0[y]).
        \]
  \item Той потърси значението на някоя дума, и нейната дефиниция не е на правилен английски, например:
        \begin{center}
          \textit{красноречив -- flooent or pursuasive in speeking or wryting}
        \end{center}
        Това съответства на глупости от вида:
        \[
          \varphi_1[x] \leftrightharpoons \exists y \; \varphi_0[y] \logand x \doteq s(y).
        \]
        Скобите имплицитно се слагат така (и въобще броя на свободни променливи не си съответства):
        \[
          (\exists y \; \varphi_0[y]) \logand x \doteq s(y).
        \]
        За това ако искаме коректна дефиниция на формула, си слагаме ръчно скобите:
        \[
          \varphi_1[x] \leftrightharpoons \exists y (\varphi_0[y] \logand x \doteq s(y)).
        \]
\end{itemize}

\section{Неуважение на кванторите $\forall$ и $\exists$}

\subsection*{Кванторът $\forall$}

Дефиницията на $\models$ ни казва, че:
\[
  \calA \models \exists x \varphi \db{b_1, \dots, b_n} \longleftrightarrow \textbf{ ЗА ВСЯКО } a \in |\calA| \text{ е изпълнено, че } \calA \models \varphi \db{a, b_1, \dots, b_n}.
\]
Семантиката на $\forall$ е да обходи \textbf{ВСЕКИ} елемент на универсума, дори когато универсума е многосортен.

Нека да разгледаме структурата $\calS = \opair{ \N \cup \calP(\N), p^\calS, q^\calS}$, където:
\begin{align*}
  p^\calS(n, m) & \dequiv n, m \in \N \text{ и } n \leq m \\ 
  q^\calS(n, X) & \dequiv n \in \N, X \subseteq \N \text{ и } n \in X.
\end{align*}
Един (грешен) опит да се определи $\{ 0 \}$ е:
\[
  \varphi_0[x] \leftrightharpoons \forall y \; p(x, y).
\]
Но тогава:
\[
\calS \models \varphi_0\db{n} \longleftrightarrow \text{ за всяко } t \in |\calS| \text{ имаме } p^\calS(n, t).
\]
В момента, в който кванторът обходи елемент $t$ на $\calP(\N)$ (което ще стане, понеже $\calP(\N) \subseteq |\calS|$), няма да бъде изпълнено $p^\calS(n, t)$.
Обикновено хората имплицитно си ограничават кванторите, но това променя дефиницията.
Дефиниции \textbf{НЕ СЕ} променят.

Правилното нещо е първо да се определи сорта с формулата:
\[
  \varphi_\N[x] \leftrightharpoons \exists y \; q(x, y).
\]
След това спокойно можем да напишем формулата:
\[
  \varphi_0[x] \leftrightharpoons \forall y (\varphi_\N[y] \Rightarrow p(x, y)).
\]
Тогава:
\[
  \calS \models \varphi_0\db{n} \longleftrightarrow \text{ за всяко } t \in |\calS| \text{ е изпълнено, че ако } t \in \N \text{, то } p^\calS(n, t).
\]
Това вече е работи както очакваме, защото експлицитно поставихме ограничение на работата на квантора.

Обикновено кванторът $\forall$ върви ръка за ръка с $\Rightarrow$, а не с $\&$.

\subsection*{Кванторът $\exists$}

Дефиницията на $\models$ ни казва, че:
\[
  \calA \models \exists x \varphi \db{b_1, \dots, b_n} \longleftrightarrow \textbf{ ИМА } a \in |\calA| \text{, за което } \calA \models \varphi \db{a, b_1, \dots, b_n}.
\]
Семантиката на $\exists$ е да намери поне един елемент от \textbf{ЦЕЛИЯ} универсум, което при многосортен универсум може да включи неща, за които не искаме да говорим.

Нашата формула може да се опитва да експлоатира, че в нашия сорт \textbf{X} няма свидетели за нещо, но е напълно възможно в целият универсум \textbf{U} това да не е така:
\begin{center}
\begin{tikzpicture}[scale=2]
  \draw[decorate,
        decoration={random steps, segment length=6pt, amplitude=2.5pt}]
    plot[smooth cycle] coordinates {
      (-1,0) (4,0.5) (5,3) (2.5,5) (-1,4) (-2,2)
    };
  \filldraw[fill=gray!30,
            draw=black,
            decorate,
            decoration={random steps, segment length=5pt, amplitude=2pt}]
    plot[smooth cycle] coordinates {
      (0.5,1.2) (3,1.5) (3.5,3.2) (2,4) (0.5,3.5) (0,2.3)
    };
  \node at (2,2.7) {\small няма свидетел тук};
  \node at (2,3.5) {\small сортът \textbf{X}};
  \fill (3.8,1.5) circle (1pt);
  \node[right] at (3.8,1.5) {\small свидетел};
  \node at (2,4.4) {\small в универсума \textbf{U}, но отвъд \textbf{X}};

\end{tikzpicture}
\end{center}
Съответно е възможно да не работи формула от вида:
\[
  \varphi[x] \leftrightharpoons \exists y \; \psi[x, y].
\]
Тогава ще трябва да я заменим с:
\[
  \varphi[x] \leftrightharpoons \exists y (\varphi_{\textbf{X}}[y] \logand \psi[x, y]).
\]
Обикновено кванторът $\exists$ върви ръка за ръка с $\&$, а не с $\Rightarrow$.

\section{Бъркане на $\varepsilon, \varnothing, \{ \varepsilon \}$ и $\{ \varnothing \}$}

Това са четири \textbf{РАЗЛИЧНИ} обекти:
\begin{itemize}
  \item $\varepsilon$ е дума (редица от символи), която има дължина $0$;
  \item $\varnothing$ е множество, което няма елементи;
  \item $\{ \varepsilon \}$ е множество, което има един единствен елемент празната дума;
  \item $\{ \varnothing \}$ е множество, което има един единствен елемент празното множество.
\end{itemize}
Обектите $\varepsilon, \{ \varepsilon \}$ и $\{ \varnothing \}$ са обекти от три \textbf{РАЗЛИЧНИ} типа:
\begin{itemize}
  \item $\varepsilon$ е \textbf{ДУМА};
  \item $\{ \varepsilon \}$ е \textbf{МНОЖЕСТВО ОТ ДУМИ};
  \item $\{ \varnothing\}$ е \textbf{МНОЖЕСТВО ОТ МНОЖЕСТВА}.
\end{itemize}
Обектите $\varnothing$ и $\varepsilon$ се различават тривиално по тип, докато останалите обекти се различават тривиално от $\varnothing$ по брой елементи.

Представете си, че сравнявате:
\begin{itemize}
  \item празен лист хартия ($\varepsilon$);
  \item празен джоб ($\varnothing$);
  \item джоб, в който има един празен лист хартия ($\{ \varepsilon \}$);
  \item джоб, в който има един празен портфейл ($\{ \varnothing \}$).
\end{itemize}

\end{document}
